\section{Full-turn integral and contact estimates}\label{app:full-turn-estimates}

\subsection{Cap area and the signed niche integral}

Green's formula expresses both cap area and niche area through their
supports. The cap calculation follows the outer boundary. The niche
calculation joins a companion tangent, a reversed corner arc, and a
first tangent. Keeping these arcs oriented makes the formulas valid
even when a parametrized graph folds. Their difference will produce
the support functional of \cref{sec:calibration}.

For vectors $z,z'\in\R^2$, write $\det(z,z')$ for their determinant.
All oriented integrals below retain their sign when endpoints are
reversed.

\begin{lemma}[Cap support area]\label{lem:cap-support-area}
Let $U$ be a downward cap of height one, with quarter supports
$F(t)=h_U(t)$ and $G(t)=h_U(t+L)$. Suppose $F,G\in C^1[0,L]$
and their first derivatives are absolutely continuous. Then
\begin{equation}\label{eq:cap-support-area}
 2|U|=\int_0^L(F^2-F'^2+G^2-G'^2)\,\dd t.
\end{equation}
\end{lemma}
\begin{proof}
The upper boundary, oriented from right to left, has support
parametrization $z(\theta)=h(\theta)u(\theta)+h'(\theta)u^\perp(\theta)$.
On regular arcs
$\det(z,z')=h(h+h'')$ almost everywhere.
At the top normal $L$, the parametrization jumps along the top
segment. Its determinant integral is its horizontal length, because
its height is one. Put its endpoints at $x=x_{\rm tl},x_{\rm tr}$.
The one-sided support derivatives are
$h'(L-)=-x_{\rm tr}$ and $h'(L+)=-x_{\rm tl}$.
Integration by parts on the two quarter arcs therefore subtracts
exactly $x_{\rm tr}-x_{\rm tl}$, cancelling the top segment's
contribution.

At the right axis normal the derivative is the right end height
$e_R$; at the left normal it is $-e_L$. The remaining boundary
terms are $l_Ue_L-r_Ue_R$. The two vertical end faces contribute
$r_Ue_R-l_Ue_L$, and the baseline contributes zero.
Green's area formula now gives \cref{eq:cap-support-area}.
The absolutely continuous integration by parts is valid without
pointwise second derivatives. Approximation by polygonal boundary
arcs, or by smooth convex caps, justifies the same Green formula
when a curvature density vanishes on an interval.
\end{proof}
For two supports with top traces $F(L)=G(0)=1$, define
\[
 \begin{gathered}
 z=(F-1)u+(G-1)u^\perp,\\
 B=(F-1)u+F'u^\perp=z+pu^\perp,\\
 D=(G-1)u^\perp-G'u=z-qu.
 \end{gathered}
\]
Then
\begin{equation}\label{eq:corner-determinant}
 \det(z,z')=(F-1)^2+(G-1)^2+(F-1)G'-(G-1)F'.
\end{equation}

\begin{lemma}[Global contact from unit curvature]\label{lem:unit-global-contact}
Suppose the actual support densities satisfy $0\le F''+F\le1$
and $0\le G''+G\le1$ almost everywhere. At every interior angle $s$,
the following points lie on the graph of the complete positive niche:
\[
 B(s)\quad\text{if }p(s)\le0,\qquad
 D(s)\quad\text{if }q(s)\ge0,\qquad
 z(s)\quad\text{if }p(s)\le0\le q(s).
\]
\end{lemma}
\begin{proof}
The complementary profile $k=1-F$ has $k''+k\ge0$.
Variation of constants, or two integrations by parts, gives
\[
 k(t)-k(s)\cos(t-s)-k'(s)\sin(t-s)
 =\int_s^t(k''+k)(r)\sin(t-r)\,\dd r\ge0
\]
for $s,t\in[0,L]$, with the integral interpreted as oriented.
Thus $F(t)-1\le\ip{B(s)}{u(t)}$ for every $t$.
The analogous inequality is
$G(t)-1\le\ip{D(s)}{u^\perp(t)}$.
In particular $B_y,D_y\ge0$, by using $t=L$ and $t=0$.

At $B(s)$ the first-wall equality holds; the second-wall height is
at least $B_y$ precisely when $p(s)\le0$.
At all other angles the first-wall height is at most $B_y$.
Taking the supremum of the two-wall minima proves its niche contact.
The proof for $D(s)$ is the same.
For a corner $z(s)$ use the identities
\[
 \ip{z(s)}{u(t)}=\ip{B(s)}{u(t)}-p(s)\sin(t-s),
\]
\[
 \ip{z(s)}{u^\perp(t)}=\ip{D(s)}{u^\perp(t)}+q(s)\sin(s-t).
\]
For $t\ge s$ the first gives the required wall alternative, and
for $t\le s$ the second does. Both equalities hold at $s$.
Also $z_y=B_y-p\cos s\ge0$. These facts prove global contact
with the whole angular supremum, not only with neighboring angles.
\end{proof}

\begin{lemma}[Signed-roof integral]\label{lem:signed-roof}
Suppose $F,G\in C^1[0,L]$ have absolutely continuous first derivatives,
$F(L)=G(0)=1$, and $0\le b\le d\le L$ satisfy $p(b)=q(d)=0$.
Let $N$ be continuous, and suppose
\[
 \begin{aligned}
 N(D_x(t))&=D_y(t)&& (0<t<d),\\
 N(z_x(t))&=z_y(t)&& (b<t<d),\\
 N(B_x(t))&=B_y(t)&& (b<t<L).
 \end{aligned}
\]
Then
\begin{equation}\label{eq:signed-roof}
 2\int_{-G'(0)}^{-F'(L)}N(x)\,\dd x
 =-\int_0^L\det(z,z')\,\dd t
         +\int_b^L p^2\,\dd t+\int_0^d q^2\,\dd t.
\end{equation}
In particular, if $p\ge0$ before $b$, $p\le0$ after $b$,
$q\ge0$ before $d$, and $q\le0$ after $d$, the last two integrals
are $\int_0^L p_-^2$ and $\int_0^L q_+^2$.
\end{lemma}
\begin{proof}
For an absolutely continuous graph parametrization $Z=(X,Y)$ with
$Y=N(X)$, the chain rule applied to $\int^X N$ gives
$\int_X^{X'}N=\int YX'$. Also
\[
 2YX'=(XY)'-\det(Z,Z').
\]
This calculation needs no monotonicity of $X$.
Join the three oriented graph arcs $D:0\to d$, $z:d\to b$,
and $B:b\to L$. The joins agree because $p(b)=q(d)=0$.
Their outer endpoints have height zero and abscissae $-G'(0)$ and
$-F'(L)$. The $XY$ terms cancel, giving
\[
 2\int N=-\int_0^d\det(D,D')
                  +\int_b^d\det(z,z')-\int_b^L\det(B,B').
\]
Direct differentiation gives
\[
 \begin{aligned}
 \det(B,B')&=\det(z,z')-p^2+[(F-1)p]',\\
 \det(D,D')&=\det(z,z')-q^2+[(G-1)q]'.
 \end{aligned}
\]
Their endpoint terms vanish by the contact and top equations.
Combining the remaining corner integrals proves
\cref{eq:signed-roof}. The sign assumptions give the final assertion.
\end{proof}

\begin{corollary}[Actual horizontal score]\label{cor:horizontal-score-identity}
Suppose a height-one horizontal canonical cap has unit curvature
on both quarters, its top face is precisely $J_U$, and there are
ordered contact switches $b,d$ with the signs in
\cref{lem:signed-roof}. Then
\[
 P_L(U)=\mathcal A(F,G)-w_U/2.
\]
\end{corollary}
\begin{proof}
By \cref{lem:unit-global-contact}, all three graph identities hold
for $N=n_L(U)$. The top face gives
$-G'(0)=l_U+w_U/4$ and $-F'(L)=r_U-w_U/4$.
Subtract \cref{eq:signed-roof} from \cref{eq:cap-support-area}
and use \cref{eq:corner-determinant}. Since $f_U=1$ on $J_U$,
the score is $|U|-w_U/2-\int_{J_U}N$.
\end{proof}

\subsection{The three-tent estimate}

The diagonal hallway bounds the exterior roof and supplies one middle
tent. The two adjacent hallways add more niche area unless their
supports move inward; moving those supports inward cuts area from the
exterior roof. We combine these two payments by completing weighted
squares. The estimate uses only three visited angles and the endpoint
laws. Put
\[
 \begin{gathered}
 r=\sqrt2,\quad k=r-1,\quad c=\cos(\pi/8),\quad s=\sin(\pi/8),\\
 A=1-k,\quad b=2-k,\quad d=3r-1,\quad \gamma=1-r/2.
 \end{gathered}
\]
These letters are local to this subsection. In particular $c,s$ do
not refer to a terminal angle, and $A$ is not the reference parameter.
Useful identities are
\begin{equation}\label{eq:three-tent-constants}
 \frac sc=k,\quad k^2=1-2k,\quad \frac1s+\frac1c=4c,
 \quad \frac r{4c^2}=k,\quad 2d+b^2=9.
\end{equation}
Let a cap have projection $[-a,a]$, middle half $[-a/2,a/2]$,
and affine middle roof rising from $1-h$ to one. Lift the left wing
by $h$ and put the middle roof at one; retain the right wing.
Concavity is preserved, because the left slopes are at least $h/a$
and the right slopes are nonpositive. Denote the lifted cap by
$\widehat U$, and define
$u=rh_{\widehat U}(\pi/4)-1$ and
$v=rh_{\widehat U}(3\pi/4)-1$.
Then $a/2\le u,v\le a$. Put $m=(u+v)/2$ and define
\begin{align}\label{eq:three-tent-objective}
 F_h(a,u,v)&=a-\frac{(a-u)^2+(a-v)^2}{2}-\frac{ah}{2}
 -(m-h/2-k)^2\notag\\
 &\hspace{12mm}+\frac{(u-k-a/2)_+^2+(v-h-k-a/2)_+^2}{2},\\
 S(a,m,h)&=a/k+(k-2)m+3/2+2k-4c+h/4,\notag\\
 G_h(a,u,v)&=F_h(a,u,v)-S(a,m,h)_+^2/d.\notag
\end{align}

\begin{lemma}[Full three-tent score]\label{lem:full-three-tent-score}
Suppose $37/25\le a\le8/5$, $0\le h\le1/3$, and the actual
endpoint laws are
\[
 e_R=\tfrac12+\tfrac h4+(3n_+-n_-)/4,\qquad
 e_L=\tfrac12-\tfrac{3h}4+(3n_--n_+)/4,
\]
where $n_\pm=n_L(U)(\pm a/2)$.
With $t=2c-r$ and $d_0=3/5-t$,
\begin{equation}\label{eq:full-three-tent-score}
 P_L(U)\le G_h(a,u,v)+
 L(h),\qquad
 L(h)=\gamma[(d_0-h/4)^2+(d_0+3h/4)^2].
\end{equation}
\end{lemma}
\begin{proof}
The first diagonal support of $U$ is unchanged by lifting, because
$(a/2,1)$ dominates all lifted left-wing points. The second diagonal
parameter is $v-h$. Thus the actual diagonal niche is
$[\min(u-k-x,v-h-k+x)]_+$.
Its apex lies in the window and is positive. Its exact integral is
\[
 (m-h/2-k)^2
  -\tfrac12[(u-k-a/2)_+^2+(v-h-k-a/2)_+^2].
\]
The lifted roof lies below
$f_0(x)=\min(1,1+u-x,1+v+x)$.
Let $D$ be its exterior deficit from this roof. Integrating $f_0$
and undoing the left lift gives
$P_L(U)\le F_h-D-(\int_J n_L-\int_Jn_{45})$.

The box-wall bound gives $n_\pm\le1-\sqrt{1-(a/2)^2}\le2/5$.
The endpoint laws then imply
\[
 e_R\ge2/5+h/4,\quad e_L\ge2/5-3h/4,
 \quad e_R+e_L\ge1-h/2.
\]
Writing $w_R=a-u$ and $w_L=a-v+h$, support containment gives
\[
 e_R\le1-w_R,\quad e_L\le1-w_L,
 \quad w_R\le3/5-h/4,\quad w_L\le3/5+3h/4.
\]
Let
$\delta_L=c+sv-h_{\widehat U}(5\pi/8)$ and
$\delta_R=c+su-h_{\widehat U}(3\pi/8)$.
Both are nonnegative. On the two flat pieces of $f_0$, the near
vertical supporting lines cut disjoint triangles of total area
$D_0=r(\delta_L^2+\delta_R^2)\le D$.
They fit before the middle window since
$\delta_L\le s(v-a/2)$ and $\delta_R\le s(u-a/2)$.

The niche at $\pi/8$ dominates the tent
\[
 \begin{gathered}
 f_R(x)=\min(b_R+kx,B_R-x/k),\\
 b_R=1-1/c+kv-h-\delta_L/c,\qquad B_R=e_R+a/k-1/s.
 \end{gathered}
\]
Against the descending diagonal branch its signed peak is
$Z_R=(B_R+b_R)/2-u+k$, with crossover endpoints
\[
 x_L=(u-k-b_R)/r,\qquad x_R=(B_R-u+k)/r.
\]
The reflected construction at $3\pi/8$ gives
\[
 \begin{gathered}
 b_L=1-1/c+ku-\delta_R/c,\quad B_L=e_L+a/k-1/s,\\
 Z_L=(B_L+b_L)/2-v+h+k.
 \end{gathered}
\]
The two triangles lie on opposite sides of the diagonal apex:
their crossover displacements from it are
$k(m\pm h/2)+2s-1+\delta/(rc)>0$.
Indeed $m-h/2\ge3/4$, $k>2/5$, and $2s>3/4$.
Their outer endpoints satisfy
$x_R\le1+a-2c<a/2$, since $a/2\le4/5$ and $c>9/10$;
the reflected endpoint satisfies the same inequality.
Thus both complete signed gain triangles fit inside $J$ whenever
their peaks are positive.

For a signed triangle of height $Z>0$ and slopes $r,-r$, whose
baseline zero exceeds its right endpoint by $\Delta$, the positive
floor loses at most $\gamma\Delta_+^2$ from the signed area
$Z^2/r$. To check this, use the coordinate $q=x_R-x$.
If $0<\Delta\le rZ$, the loss is
\[
 \int_0^{k\Delta}rq\,\dd q+
 \int_{k\Delta}^{\Delta}(\Delta-q)\,\dd q
   =\tfrac12[rk^2+(1-k)^2]\Delta^2=\gamma\Delta^2.
\]
If $\Delta>rZ$, the remaining positive tent has area
$r[(1+1/r)Z-\Delta]_+^2$.
When its bracket is nonnegative, subtracting
$Z^2/r-\gamma\Delta^2$ leaves
$(1+r/2)(\Delta-rZ)^2\ge0$.
When the bracket is negative the latter bound is already
nonpositive. The case $Z\le0$ needs no positive interval.
For the two triangles the losses satisfy
$\Delta_R\le(w_R-t)_+$ and $\Delta_L\le(w_L-t)_+$,
so their total is at most $L(h)$.

Write $Z_R=Z_R^0-\delta_L/(2c)$ and
$Z_L=Z_L^0-\delta_R/(2c)$.
The endpoint-height sum gives $Z_R^0+Z_L^0\ge S(a,m,h)$.
If $H_0=((Z_R)_+^2+(Z_L)_+^2)/r$, Cauchy--Schwarz gives
\[
 S_+\le\sqrt{2rH_0}+\sqrt{rD_0/(4c^2)},\qquad
 S_+^2\le(2r+r/(4c^2))(H_0+D_0)=d(H_0+D_0).
\]
Combining the exterior deficit, diagonal niche, additional two
triangles, and floor correction proves
\cref{eq:full-three-tent-score}. Each exterior deficit is charged
once, and the two niche gain regions are disjoint.
\end{proof}

\begin{lemma}[Stationarity and the width derivative]\label{lem:three-tent-general-plane}
Fix $h\ge0$ and a comparison width $a_0$. Put
\[
 \sigma=s_0(a_0)-Ah/4,\quad
 s_0(a)=da/2-(4c-2),\quad
 m_*=(a_0+k)/2+b\sigma/9+h/4.
\]
Suppose
\[
 \sigma\ge0,\qquad m_*\le k+a_0/2,\qquad
 r-a_0-2d\sigma/9\le0.
\]
Then for $a\ge a_0$ and diagonal data with nonnegative apex,
\[
 G_h(a,u,v)\le
 B(a_0)-kh/2-h^2/8-2\sigma^2/9,
 \qquad B(a)=a-(a-k)^2/2.
\]
\end{lemma}
\begin{proof}
Use coordinates $(a,m,\Delta)$, with $\Delta=(u-v)/2$.
Before the two clipping terms the Hessian of $F_h$ is
\[
 \begin{pmatrix}-2&2&0\\2&-4&0\\0&0&-2\end{pmatrix}.
\]
Each active clip adds the outer product of $(-1/2,1,1)$ or
$(-1/2,1,-1)$. With both present the Hessian is
\[
 \begin{pmatrix}-3/2&1&0\\1&-2&0\\0&0&0\end{pmatrix},
\]
which is negative semidefinite. Any other clipping pattern subtracts
a positive semidefinite matrix. The gradients match at the interfaces,
and subtracting $S_+^2/d$ preserves joint concavity.

At $u=v=m$ in the unclipped region put $m'=m-h/4$ and
$m_0(a)=(a+k)/2$. Completion of squares gives
\[
 G_h=B(a)-kh/2-h^2/8-2(m'-m_0(a))^2
       -[s_0(a)-b(m'-m_0(a))-Ah/4]^2/d.
\]
Since $2d+b^2=9$, both support derivatives vanish at
$(a,u,v)=(a_0,m_*,m_*)$. The clipping hypothesis makes the
right clip zero; the left clip is smaller by $h$. Also
$S=2d\sigma/9\ge0$. The value there is the asserted bound,
and its width derivative is $r-a_0-2d\sigma/9$.
The tangent plane of a concave function is an upper bound everywhere.
Its support derivatives vanish and its width derivative is nonpositive,
so it bounds every larger width by this value.
\end{proof}

\begin{lemma}[Three-tent supporting plane]\label{lem:full-three-tent-critical}
Fix $0\le h\le1/3$ and put
\[
 \begin{gathered}
 a_0=37/25,\qquad B(a)=a-(a-k)^2/2,\\
 s_0(a)=da/2-(4c-2),\qquad m_0(a)=(a+k)/2.
 \end{gathered}
\]
On $a\ge a_0$, $a/2\le u,v\le a$, the function $G_h$ is
jointly concave and
\begin{equation}\label{eq:three-tent-supporting-plane}
 G_h(a,u,v)\le
 K_0(h):=B(a_0)-kh/2-h^2/8
             -\tfrac29[s_0(a_0)-Ah/4]^2.
\end{equation}
\end{lemma}
\begin{proof}
Apply \cref{lem:three-tent-general-plane}. At $a_0=37/25$ its
stationary support coordinate is
\[
 u_*=v_*=(a_0+k)/2+\frac b9[s_0(a_0)-Ah/4]+h/4.
\]
To check clipping, keep the angle parameters linked by
$k^2+2k=1$ and $2c^2(1-k)=1$. With
$\mathcal N_q(x)=(x+q/x)/2$, define
\begin{equation}\label{eq:triangle-newton-parameters}
 \begin{gathered}
 k_+=\mathcal N_2(\mathcal N_2(3/2))-1,
 \qquad k_-=2/(1+k_+)-1,\\
 Q_\pm=2-2k_\pm,\qquad
 c_-=[\mathcal N_{Q_-}(\mathcal N_{Q_-}(1))]^{-1},\qquad
 c_+=\mathcal N_{Q_+}(\mathcal N_{Q_+}(1))/Q_+.
 \end{gathered}
\end{equation}
Squaring the Newton inequality gives $k_-<k<k_+$ and
$c_-<c<c_+$. Write
$s_+=(2+3k_+)a_0/2-(4c_--2)$.
Substitution in these formulas gives $2/3<s_0(a_0)<s_+<3/4$
and $s_0(a_0)-Ah/4>0$. Since $bA=3-5k$, the larger clip obeys
\[
 u_*-k-a_0/2
 <-k_-/2+(2-k_-)s_+/9+(6+5k_+)(1/3)/36<0.
\]
The last inequality follows by inserting the Newton formulas;
no separately chosen bound on the support coordinate is needed.
The other clip is smaller by $h$.
Also $a_0/2<u_*<a_0$, and $S>0$ there.
The width derivative is negative because
\[
 r-a_0-\frac{2d}{9}[s_0(a_0)-Ah/4]<r-a_0<0.
\]
These checks verify the hypotheses of
\cref{lem:three-tent-general-plane}, which gives the bound $K_0(h)$
for all $a\ge a_0$.
\end{proof}

The floor corrections in the full, proper, and dry-endpoint comparisons
have the same form. The following payment separates the geometric loss
from the final arithmetic check.

\begin{lemma}[Tilt pays for the floor]\label{lem:triangle-tilt-payment}
Let $\sigma=s_0(a_0)$ and put
\[
 \Psi(h)=B(a_0)-kh/2-h^2/8-\tfrac29(\sigma-Ah/4)^2
       +\gamma[(p+\lambda h)^2+(q+\mu h)^2].
\]
If $k\ge2/5$, $k\le1$, $\sigma\le3/4$,
$\lambda p+\mu q\le1/6$, $\lambda^2+\mu^2\le5/8$, and
$0\le h_0\le h$ with $h+h_0\le1/2$, then $\Psi(h)\le\Psi(h_0)$.
\end{lemma}
\begin{proof}
Subtract at $h_0$ and factor out $h-h_0$. The remaining factor is
\[
 -k/2+A\sigma/9+A(\lambda p+\mu q)
 +(h+h_0)[-1/8-A^2/72+A(\lambda^2+\mu^2)/2].
\]
Since $0\le A\le3/5$, it is at most
\[
 -1/5+1/20+1/10+1/32<0.
\]
The first term is the direct tilt cost; the next two pay for
moving the stationary point and the floor cuts. The last pays
for their quadratic remainder.
\end{proof}

For the three applications, $t=2c-(1+k)$ and the data are
\[
\begin{array}{c|cccc|c}
 &p&q&\lambda&\mu&h_0\\ \hline
 \text{full}&3/5-t&3/5-t&-1/4&3/4&0\\
 \text{proper}&0&2/3-t&0&2/3&0\\
 \text{dry endpoint}&1/2-t&7/12-t&-1/4&3/4&0
\end{array}
\]
The Newton enclosures give $5/12<t<1/2$.
Thus $\lambda p+\mu q\le1/6$ in every row, and the largest
$\lambda^2+\mu^2$ is $5/8$, from the endpoint-law slopes
$-1/4,3/4$. The allowed heights satisfy $h+h_0<1/2$.
One tilt payment therefore covers all three exclusions.

\paragraph{The width fixed by the initial radius}

A dry low endpoint makes the initial momentum
$B=b(h)+3n_+/4$, where $b(h)=-1/2+5h/4$.
The radius test allows span $R(h)=\sqrt{5-b(h)^2}$.
The box gives $n_+\le1/3$, so the largest momentum at height $h$
is $b(h)+1/4$. Choose the height where its floor-energy test meets
four at that allowed radius:
\begin{equation}\label{eq:tilted-balance-height}
 R(h_0)^2+[b(h_0)+1/4]^2+[b(h_0)+1/4](1-h_0)
 -R(h_0)\sqrt{h_0(2-h_0)}=4.
\end{equation}
This determines a unique $h_0\in(0,1/3)$.
To check existence and obtain bounds, put
\[
 \begin{aligned}
 V(h)&=R(h)^2+[b(h)+1/4]^2+[b(h)+1/4](1-h)-4,\\
 F(h)&=V(h)^2-R(h)^2h(2-h).
 \end{aligned}
\]
Here $V$ is positive; $F$ decreases on $[0,1/3]$ and changes sign
between its endpoints. The squared equation therefore gives the
same height. Starting with $[l,u]=[0,1/3]$, three simultaneous updates
\[
 (l,u)\longmapsto\left(l-\frac{F(l)}{F'(l)},\,
                        u-\frac{F(u)}{F'(l)}\right)
\]
give bounds $h_-\le h_0\le h_+$.
The divided difference of $F$ lies above $F'(l)$ on each bracket,
so these are enclosing updates, rather than fitted approximations.
The shared radius width is now
\begin{equation}\label{eq:tilted-radius-width}
\begin{gathered}
b_0=b(h_0),\quad R_0=R(h_0),\quad C_{\rm rad}=R_0/3,\\
R_+=\mathcal N_{R(h_+)^2}^{\circ2}(2),\qquad
R_-=R(h_-)^2/R_+,\qquad a_{\rm rad}=2R_-/3.
\end{gathered}
\end{equation}
The span lower bound $d\ge2$ supplies the Newton seed.
The brackets give $R_-<R_0<R_+$.
The triangle comparison below excludes every larger width, even when
$h<h_0$, and can be evaluated at zero tilt like the other two rows.
Both boundaries are therefore fixed by the two energy tests.

\begin{lemma}[A common endpoint comparison]\label{lem:triangle-newton-corners}
Use $k_+,c_+$ from \eqref{eq:triangle-newton-parameters} and define
\[
\begin{aligned}
 E&=(a+1-h/2)k+a-(a^2+1)/2-h^2/8\\
  &\quad-2z^2/9+\tfrac A2(wD_R^2+D_L^2),\\
 z&=(2+3k)a/2-(4c-2)-Ah/4,\\
 D_i&=b_i-(2c-1-k).
\end{aligned}
\]
For each of the following geometric endpoint data,
$E(a,h,k,c;b_R,b_L,w)\le E(a,h,k_+,c_+;b_R,b_L,w)<\underline p$:
\[
\begin{array}{c|ccccc}
 &a&h&b_R&b_L&w\\ \hline
 \text{full}&37/25&0&3/5&3/5&1\\
 \text{proper}&37/25&0&2/3&2/3&0\\
 \text{dry endpoint}&a_{\rm rad}&0&1/2&7/12&1
\end{array}
\]
Here $w=0$ records the proper comparison's single joint floor square.
\end{lemma}
\begin{proof}
The identity $k^2+2k=1$ rewrites $B(a)$ as the linear term in
$E$. To bound the whole expression, first increase $c$ to $c_+$,
then increase $k$ to $k_+$. On these segments, the Newton enclosures
give $3/5\le z\le3/4$ and $0\le D_i\le1/4$.
The two changes both raise $E$:
\[
 \partial_cE=16z/9-2A(wD_R+D_L)>0,
\]
\[
\begin{aligned}
 \partial_kE&=a+1-h/2-\tfrac49(3a/2+h/4)z\\
 &\quad-\tfrac12(wD_R^2+D_L^2)+A(wD_R+D_L)\\
 &\ge a/2+15/16-7h/12>0.
\end{aligned}
\]
Here $0\le w\le1$, and all three rows have $h=0$.
The first comparison retains the improvement in the stationary square
alongside the loss in floor payment; the second retains the linear
area reward. Both can equally be obtained by subtracting the two
squares. Inserting the same Newton formulas in each row gives
$E<\underline p$ using \eqref{eq:reference-half-area-lower}.
\end{proof}

\begin{corollary}[Full width cut]\label{cor:full-three-tent-cut}
Every cap satisfying the hypotheses of
\cref{lem:full-three-tent-score} has $P_L(U)<\underline p$.
\end{corollary}
\begin{proof}
By \cref{lem:full-three-tent-critical},
$P_L(U)\le K_0(h)+L(h)$.
The full row of \cref{lem:triangle-tilt-payment} bounds this at $h=0$.
Its value is the full row of \cref{lem:triangle-newton-corners},
which is below $\underline p$.
\end{proof}

\subsection{The first-wing energy test}
\label{app:first-wing-energy}

The following estimates use the tilted notation of
\cref{sec:full-turn}. In particular $I=[-2C,2C]$, $J=[-C,C]$,
the middle roof rises from $1-h$ to $1$, and $F,G$ are the
high-wing and low-wing supports. Write
\[
 d=3C,\qquad B=e_R-1+h,\qquad
 s_h=\sqrt{h(2-h)},\qquad H(C)=1-\sqrt{1-C^2}.
\]

The target $u\le1$ follows if $q\le1$ whenever the first source can
be exposed. At its first $p$-zero this means $q+1\le2$. Before the
corner reaches the floor, the support trajectory is a free rotation.
The two tests below bound its energy depending on which event comes
first; the constants four and five are the corresponding squared
radii.

\begin{lemma}[Initial-floor test]\label{lem:tilted-initial-energy}
Suppose $1/2<C<4/5$ and $0<h<1/3$. At a canonical full
maximum, the two inequalities
\begin{equation}\label{eq:tilted-initial-energy}
 d^2+B^2\le5,\qquad
 d^2+B^2+B(1-h)-d s_h\le4
\end{equation}
imply $u\le1$ almost everywhere on $(0,L)$.
\end{lemma}
\begin{proof}
On the initial component where $p,q>0$ and $D_y<0$, both source
densities vanish. The first vanishes by same-sign shadowing.
A positive second-source contact with companion slack would have
to be the stationary second-wall contact $D$, whose height is
negative. A corner is shadowed because both its walls increase
with the angle. The source identity therefore makes the second
density zero as well. This reasoning also excludes a floor
contact: a nonstationary one with slack is raised by a nearby
angle, and a same-sign floor corner is raised on both walls.
Compact exhaustion in angle justifies the statement at the
level of the limiting source measures.

Consequently, until this component meets the floor or $p=0$,
\begin{equation}\label{eq:initial-floor-solution}
 \begin{gathered}
 F=2C\cos t+e_R\sin t,\qquad
 G=C\sin t+(1-h)\cos t,\\
 D_y=1-h-\cos t,\\
 p-1=B\cos t-d\sin t,\qquad
 q+1=d\cos t+B\sin t.
 \end{gathered}
\end{equation}
Put $\gamma_0=\arccos(1-h)$ and let $\alpha$ be the first
$p$-zero on the initial positive-$q$ component. Stop at
$\min(\alpha,\gamma_0)$. The rotating vector in
\cref{eq:initial-floor-solution} has squared norm $R^2=d^2+B^2$.
Thus
\[
 \begin{aligned}
 E_{++}:=(p-1/2)^2+(q+1)^2&=R^2+p-3/4\\
 &\le R^2+1/4+\max\{B(1-h)-ds_h,-1\}
 \end{aligned}
\]
at the stopping time. The $-1$ alternative is the value when
$p=0$; the other alternative is the value at the floor.
Both tests in \cref{eq:tilted-initial-energy} say that this upper
bound is at most $17/4$.

If the floor is reached first, the same-sign bounds $u=0$,
$v\le1/2$ give
\[
 E_{++}'=2(p-1/2)(-1-q)+2(q+1)(v-1+p)
        =2(q+1)(v-1/2)\le0.
\]
At the first $p$-zero, $1/4+(q+1)^2\le17/4$, hence $q\le1$.
While $p<0$ and $q>0$, the arm estimate gives $q'\le-1/2$:
for $-1\le p\le0$ use $v\le(1-p)/2$, and for $p<-1$
use $v\le-p$. The positive-$q$ component therefore has $q\le1$
after this zero.

Any later positive-$q$ component starts at $q=0$ with $0<p\le1$;
entry with $p\le0$ is excluded by the preceding derivative bound.
While $p,q>0$, $u=0$, $p'\le-1$ and $q'\le p-1/2$.
Its maximum $q$ is at most
$\int_{1/2}^1(x-1/2)\,\dd x=1/8$. Once $p$ is negative,
$q$ decreases again. If the initial positive-$q$ component ends
before the stopping time, only these later components need be
considered. Finally $q\ge-1$ everywhere, and
$u\le\kappa(q)$, with $u=0$ in the positive-positive sector.
The proved bounds $-1\le q\le1$ on all remaining sectors give
$u\le1$.
\end{proof}

\begin{lemma}[Short widths and the radius width]\label{lem:tilted-unit-regions}
The first density of a canonical full maximum is at most one if
\(C\le2/3\). If its low middle endpoint is dry, the same conclusion
holds for \(C\le C_{\rm rad}\) and \(h_0\le h<1/3\), with
\(C_{\rm rad},h_0\) from \eqref{eq:tilted-radius-width}.
\end{lemma}
\begin{proof}
For \(C\le2/3\), write \(e=e_R\in[0,1]\).
Then \(d^2\le4\), \(|B|=|e-1+h|\le1\) and \(s_h\ge h\).
The first test follows at once. The second follows from
\[
B^2+B(1-h)=e(e-1+h),\qquad
 e(e-1+h)-ds_h\le h(e-d)\le0,
\]
since \(d>3/2\ge e\).

Now assume \(2/3<C\le C_{\rm rad}\), \(h\ge h_0\), and \(n_-=0\).
Since \(9C_{\rm rad}^2=5-b_0^2<5\), the box height obeys
\(H(C)\le1/3\). The dry endpoint pressure therefore gives
\[
-1/2+5h/4\le B\le-1/4+5h/4.
\]
On \([h_0,1/3]\) both endpoints have absolute value at most
\(-b_0\), and hence
\[
d^2+B^2\le9C_{\rm rad}^2+b_0^2=5.
\]
This is exactly why the radius width was chosen.

For the floor test put
\(Q(d,B,h)=d^2+B^2+B(1-h)-d\sqrt{h(2-h)}\).
It increases with \(B\) on this interval, since
\(2B+1-h\ge3h/2\ge0\), and with \(d\ge2\).
Thus it suffices to use $d=R_0$ and $B=-1/4+5h/4$.
The resulting function is convex in $h$: its polynomial part
has quadratic coefficient $5/16$, and the negative semicircle is convex.
Only the two endpoints $h=h_0,1/3$ remain.
At $h_0$ the test equals four by its defining equation.
At $1/3$, use $R_0^2\le5$, $R_0\ge2$ and
$\sqrt{(1/3)(2-1/3)}\ge2/3$. These give
\[
 Q(R_0,1/6,1/3)\le5+(1/6)^2+(1/6)(2/3)-2(2/3)<4.
\]
Convexity supplies the whole interval.
Apply \cref{lem:tilted-initial-energy}.
\end{proof}

\subsection{Ordered contact geometry}

Here $F,G,u,v,p,q,z,B,D,C,h,T$ have the tilted meanings in
\cref{eq:tilted-wing-supports,eq:tilted-contact-data}. In
particular $G(0)=1-h$; it is the low-wing support, not the
support containing the middle facet.

\begin{lemma}[Contact signs]\label{lem:tilted-contact-signs}
If $1/2<C<4/5$, $0<h<1/2$ and $u\le1$ almost everywhere,
there are $0<\alpha<\eta<L$ such that
\[
 \begin{array}{c|ccc}
 t&(0,\alpha)&(\alpha,\eta)&(\eta,L)\\ \hline
 p&>0&<0&<0\\
 q&>0&>0&<0.
 \end{array}
\]
Moreover $p(\alpha)=q(\eta)=0$.
\end{lemma}
\begin{proof}
Initially $p(0)=e_R+h>0$, $q(0)=3C-1>0$, whereas
$p(L)=1-3C-T<0$ and $q(L)=-e_L<0$.
Before $\delta=\arctan(h/(2C))$, the zero-density solution
gives $q'\ge-1$ while $p$ is positive. Since
$\delta<h<1/2<q(0)$, $q$ cannot reach zero there.
After $\delta$, the arm bound gives $p\le1$.
On the positive-positive sector, $u=0$, and
\[
 [(p-1)^2+(q+1)^2]'=2(q+1)v\ge0.
\]
Its initial value is at least $9C^2>2$. A first exit through
$q=0$ would instead have $0\le p\le1$ and energy at most two.
Thus the first exit is $p=0$, at a time $\alpha$, with $q>0$.

For $q\ge0$, the arm bound gives
$p'=u-1-q\le\kappa(q)-1-q\le-1/2$.
For $p\le0$, the companion bound similarly gives $q'\le-1/2$.
Consequently $p$ stays negative while $q>0$, and $q$ reaches
zero at some $\eta>\alpha$ with $p(\eta)<0$.
In the ensuing negative-negative sector $v=0$, $u\le1/2$ and
\[
 [(p-1)^2+(q+1/2)^2]'=2(p-1)(u-1/2)\ge0.
\]
At $\eta$ this energy is strictly greater than $5/4$.
A return through $p=0$ would have $-1\le q\le0$ and energy
at most $5/4$, since the lower arm bound is $q\ge-1$.
Such a return is impossible. As long as $p<0$, the preceding
$q'\le-1/2$ bound also prohibits a return through $q=0$.
This proves all the signs, including strictness at the ends of
the open sectors. The argument uses first crossing times of
continuous functions and integrated almost-everywhere derivative
bounds, so it covers zero-density intervals.
\end{proof}

\begin{lemma}[First-unit contact geometry]
\label{lem:tilted-first-unit-geometry}
Under the hypotheses of \cref{lem:tilted-contact-signs},
$B_x$ is nondecreasing and $B_y$ is nonincreasing. Let $\beta$
be its first crossing of $C$ (use $\beta=L$ when $T=0$).
Then
\begin{equation}\label{eq:tilted-crossing-box}
 \alpha<\beta,\qquad \beta\ge\arccos C,\qquad
 \sin\beta\ge3/5,\qquad \cot\beta\le4/3.
\end{equation}
The positive niche is precisely $(x_*,C+T)$ for one
$x_*<1-2C$, and
\begin{equation}\label{eq:tilted-no-ghost-projection}
 T=(C+x_*)_+.
\end{equation}
If the companion density exceeds one, its first bad entry
$\tau$ satisfies
\begin{equation}\label{eq:tilted-first-bad-bound}
 \alpha<\beta<\tau<\eta,\quad
 p(\tau)=-1,\quad 0<q(\tau)\le1/4,
 \qquad v\le1\text{ a.e. on }(0,\tau).
\end{equation}
In particular $T=0$ implies $v\le1$ almost everywhere.
\end{lemma}
\begin{proof}
The monotonicity follows from $B'=(u-1)u^\perp$ and $u\le1$.
If $B_x\ge C$, its associated outer point
$(X,Y)=B+u(t)$ lies in the cap, so $\cos t\le C$ and
$X\ge C+\cos t$, $Y\le1$. The low endpoint gives
$G\ge C\sin t+(1-h)\cos t$, whence
\[
 p\le1-(2C+\cos t)\sin t+h\cos t<0.
\]
For the last sign put $k=\cos t\in[0,C]$. The upper bound
$1-(2C+k)\sqrt{1-k^2}+k/2$ is convex in $k$.
At zero it is $1-2C<0$; at $C$ it is negative because
\[
 9C^2(1-C^2)-(1+C/2)^2>0\quad(1/2\le C\le4/5).
\]
This polynomial has derivative
$(2-3C)(24C^2+16C-1)/2$, hence its minimum is at an
endpoint; its endpoint values are $1/8$ and $71/625$.
Thus $B_x(\alpha)<C$, proving the first inequality in
\cref{eq:tilted-crossing-box}; the outer point at the crossing
gives the remaining three.

For the positive set write
\[
 \begin{gathered}
 R_t(x)=\frac{F(t)-1-x\cos t}{\sin t},\qquad
 S_t(x)=\frac{G(t)-1+x\sin t}{\cos t},\\
 L_0(t)=\frac{1-G(t)}{\sin t}.
 \end{gathered}
\]
On $x<1-2C$, the envelope $\sup_t S_t(x)$ is finite, convex,
and nondecreasing in $x$. Its endpoint limits are $-h$ at
zero and $-\infty$ at $L$. It is at most $-h$ for
$x\le-2C$, by the low-wing box $[-2C,-C]\times[0,1-h]$;
at $x=1-2C$ its $L$-limit is $e_L>0$.
There is exactly one zero $x_*\in(-2C,1-2C)$.
Indeed a zero is attained at an interior angle and that wall
has positive slope $\tan t$, so all points to its right have
positive envelope. All points to its left have negative envelope.
Equivalently $x_*=\inf_t L_0(t)$.

The derivative of $(F-1)/\cos t$ is $B_y/\cos^2t\ge0$.
Its intercept therefore rises from $2C-1$ to $C+T$.
Every first wall is positive for $x\le1-2C<2C-1$, and
every first wall is nonpositive for $x\ge C+T$.
For $1-2C<x<C+T$, angles approaching $L$ give
$R_t(x)=(C+T-x)(L-t)+o(L-t)>0$ and $S_t(x)\to+\infty$.
This proves the positive interval. The additional actual
high-point wall is nonpositive on $J$ and does not affect it there.

The finite positive graph has horizontal length equal to the
sine/cosine moment of its two source measures. Source convergence
and \cref{eq:tilted-wing-traces} give its limiting length as
\[
 \int_0^L(u\sin t+v\cos t)\,\dd t=2C-T.
\]
Uniform niche convergence gives the lower bound from the positive
interval. To get the reverse bound, take a compact interval
strictly left of $x_*$. All surrogate second walls there are strictly
negative when their angles lie in $[\varepsilon,L-\varepsilon]$.
The additional actual high-point wall is also strictly negative
there. The finite positive graph on this compact interval must
therefore use endpoint-angle sources. Their total mass is
$O(\varepsilon)+o(1)$ by \cref{eq:variation-coarse}.
Exhaust the zero interval and then let $\varepsilon\downarrow0$.
Thus the finite positive lengths converge to
$C-\max\{-C,x_*\}$. Equating the two lengths proves
\cref{eq:tilted-no-ghost-projection}.

Let $\tau$ be the first entry into $\{p<-1,q>0\}$.
Before it, $v\le1$: use $v\le\kappa(p)$ when $p\ge-1$,
and $v=0$ in the negative-negative sector when $p<-1$.
On a level set $q=0$, its derivative is zero almost everywhere,
whereas $p<-1$ would give $q'\le-1$; no exceptional mass remains.
Entry through $q=0$ is impossible because $q'\le-1/2$ for
$p\le0$. Thus $p(\tau)=-1$, $q(\tau)>0$.

For $\alpha<t<\tau$, past $D_x$ is nondecreasing, whereas
future $B_x$ is globally nondecreasing. Hence the positive core
corner is the unique active angle: earlier second walls and
later first walls are strictly below it. It lies in $J$, since
$-C<D_x+q\cos t=z_x=B_x+p\sin t<C$.
It contributes first-source density $q$. Before $\beta$,
the disjoint visible $B$ tangency contributes another $1-u$.
The source identity gives
$u\ge(1+q)/2$ before $\beta$, and $u\ge q$ after it.
Plateaus contribute zero tangent length and do not change these
bounds. If $\tau\le\beta$, the bounds $u\le1$ and
$u\le\kappa(q)$ make $q\le1$, $u=(1+q)/2$ on the core
up to $\tau$. There
$E=(p-1)^2+(q+1)^2$ is nonincreasing and $E(\alpha)\le5$;
but $E(\tau)=4+(q(\tau)+1)^2>5$, a contradiction.
The same strict comparison at the crossing excludes equality
$\tau=\beta$.

On $(\beta,\tau)$, put $K=q+(p-1)^2/4$.
Since $u\ge q$ and $-1<p<0$, $v\le(1-p)/2$ gives
$K'\le v+(p-1)/2\le0$.
At $\beta$, $E\le5$ yields
\[
 K(\beta)\le1+q(\beta)/2-q(\beta)^2/4\le5/4.
\]
As $K(\tau)=1+q(\tau)$, this proves
\cref{eq:tilted-first-bad-bound}. If $T=0$, $\beta=L$ and
the first bad entry is impossible; the arm and same-sign
bounds then give $v\le1$ everywhere.
\end{proof}

\subsection{The joint occupation of a corner}

A visible corner belongs to both wall families. Resolving its tangent
vector into those two wall directions fixes the ratio of their source
lengths. Thus both families must use the same occupied fraction of a
corner. This shared fraction is what lets the energy argument pay for
folds without counting a contact twice.

\begin{lemma}[Common source occupation]
\label{lem:variation-joint-occupation}
Under the first-unit hypotheses, suppose the companion density
exceeds one on a set of positive measure, and let the
ordered contacts and $\beta$ be as above. There are measurable
$d,\chi\in[0,1]$ such that, almost everywhere,
\begin{equation}\label{eq:joint-source-equations}
 \begin{array}{c|cc}
 &u&v\\ \hline
 (0,\alpha)&0&(1-v)d\\
 (\alpha,\eta)&(1-u)b+q\chi&(1-v)d-p\chi\\
 (\eta,L)&0&0,
 \end{array}
 \qquad b=\mathbf1_{\{t<\beta\}}.
\end{equation}
The second tangent occupation satisfies $d=0$ where $v>1$ or
$D_y<0$. Before a first bad entry $\tau$, one has $\chi=1$
on $(\alpha,\tau)$.
\end{lemma}
\begin{proof}
Keep the source angle and graph abscissa together in the finite
polygons. Assign to every positive first-source segment its
arclength measure at $(t,x)$; call this $\mu_{i,F}$, and
define $\mu_{i,G}$ similarly. Their exact graph-vector identities
on the window interior are
\begin{align*}
 \pi_{x\#}(\sin t\,\mu_{i,F}+\cos t\,\mu_{i,G})
   &=\mathbf1_{\{N_i>0\}\cap J_i}\,\dd x,\\
 \pi_{x\#}(-\cos t\,\mu_{i,F}+\sin t\,\mu_{i,G})
   &=\dd(N_i|_{J_i}).
\end{align*}
The second is a distributional identity. Bounded mass gives a
joint weak subsequence. The no-ghost argument in
\cref{lem:tilted-first-unit-geometry} passes the first identity
to Lebesgue measure on $(x_*,C)$. Uniform graph convergence
passes the second to $N'(x)\,\dd x$ on this positive interval.
The angular marginals are exactly $u\,\dd t,v\,\dd t$ by
\cref{lem:variation-sources}. There are no endpoint-angle atoms.
The graph is locally Lipschitz on its positive interval:
on a compact positive part, the endpoint-angle walls tend to
zero or below, so all active angles remain in a compact
subinterval of $(0,L)$. These identities also exclude extra
mass at a spatial singleton or at a floor point.

At a regular positive active source with first-wall slack,
angular stationarity gives $x=B_x(t)$ and height $B_y(t)$;
this needs $p<0$. The analogous second-wall contact is $D$,
with $q>0$. A simultaneous contact is $z$.
A corner with two positive momenta is shadowed by a slightly
larger angle, and one with two negative momenta by a slightly
smaller angle. Thus only core corners remain.
At a second tangent
\[
 \partial_t^2 S_t(D_x(t))=-(1-v(t))/\cos t.
\]
When $v>1$, it is an angular local minimum with strict companion
slack, and cannot be an active full-envelope maximum. The same
conclusion follows at density points from the integrated
second-derivative inequality. Second-tangent occupation vanishes
there and when $D_y<0$.

The first-unit bound makes $B$ the global first-wall envelope.
On $(\alpha,\beta)$ its companion has strict slack. Its graph
has negative slope and lies to the right of every core corner,
because $z_x$ strictly decreases from $B_x(\alpha)$.
An active $D$ tangent has positive slope and cannot coincide on
a set of positive graph measure. Thus the entire first-tangent
contribution is $(1-u)\mathbf1_{(\alpha,\beta)}\,\dd t$.
Plateaus have zero spatial projection and hence zero extra mass.

On the measurable corner contact set $N(z_x(t))=z_y(t)$,
the chain rule and the derivative-zero property on level sets give
\[
 N'(z_x(t))=\frac{p\sin t+q\cos t}{p\cos t-q\sin t}
 \quad\text{a.e.}
\]
On a compact core interval $z_x$ is bi-Lipschitz, so null sets
in angle and abscissa cause no difficulty. Disintegrate the two
vector identities at $x=z_x(t)$. Any competing $D$ contact has
the same tangent $(1,N'(x))$; subtract its parallel flux.
For the remaining first and second arclength densities $A_F,A_G$
the vertical identity becomes
\[
 (-\cos t-N'\sin t)A_F+(\sin t-N'\cos t)A_G=0.
\]
Substitution gives $A_G=(-p/q)A_F$.
Their nonnegative horizontal share is at most one. Since the
increasing-abscissa corner displacement is $-z'(t)\,\dd t$,
one common number $\chi\in[0,1]$ therefore assigns angular
densities $q\chi$ and $-p\chi$. Exhaust compact core intervals
to obtain a measurable $\chi$ on the whole core. Strict unique
corner visibility before $\tau$ makes $\chi=1$ there.

For a visible regular second tangent, $D_x'=(1-v)\cos t$.
The area formula pulls its horizontal share back to angle as
$(1-v)d\,\dd t$, for a measurable $d\in[0,1]$.
At almost every graph abscissa there is at most one visible
second tangent, since its slope $\tan t$ determines its angle.
Set $d=0$ on absent or forbidden contacts.
All positive second tangents and core corners lie in $J$:
their abscissa exceeds $L_0(t)\ge x_*>-C$, whereas
$D_x\le1-C<C$ and $z_x\le B_x(\alpha)<C$.

These three classes exhaust the limiting source measure.
There is no nonnegative remainder, since the weighted spatial
projection already equals the whole positive graph projection.
Equating angular marginals gives
\cref{eq:joint-source-equations}. In the final sector the first
tangencies are outside $J$, as $\beta<\eta$ whenever this
occupation is used in the nonunit branch; the second density
already vanishes by same-sign shadowing.
\end{proof}

\subsection{The two clipped strip moments}

The top overhang $T$ appears as a horizontal moment at the low-wing
floor crossing and as a vertical moment at the high-wing window
crossing. Integrating the contact energy compares these moments.
Endpoint stationarity cancels the end heights, leaving the required
term $6CT$ and convolution remainders that the two unit-curvature
bounds can control.

\begin{lemma}[Two unit wings]\label{lem:tilted-two-unit}
A canonical full maximum with positive middle slope,
$1/2<C<4/5$ and $0<h<1/2$ cannot have both $u,v\le1$
almost everywhere.
\end{lemma}
\begin{proof}
Both $B$ and $D$ have nondecreasing abscissae; $B_y$ decreases
to zero and $D_y$ increases from $-h$ to $e_L>0$.
Moreover
$B_x-D_x=q\cos t-p\sin t\ge4C-2>0$.
Let $\gamma$ be a floor crossing of $D$, and $\beta$ the
right-window crossing of $B$. Ordered signs and the exact wall
derivatives
\[
 \partial_tR_t(x)=\frac{x-B_x(t)}{\sin^2t},\qquad
 \partial_tS_t(x)=\frac{x-D_x(t)}{\cos^2t}
\]
show that the positive graph in the window consists of
$D[\gamma,\eta]$, the reversed core $z[\alpha,\eta]$,
and $B[\alpha,\beta]$. Indeed at a tangent its sole active
wall is the global angular maximum with companion slack;
at a core corner every earlier second wall and later first
wall is strictly below it. The source-box estimate gives
$\alpha<\beta$, and positivity of $z_y(\eta)$ gives
$\gamma<\eta$. No relative order of $\beta$ and $\gamma$
is needed. Flat tangent parameter intervals have zero projection.

The no-ghost identity now says
$T=C+D_x(\gamma)>0$. Strict positivity also follows directly
from $h=\int_0^\gamma(1-v)\sin t>0$.
Local finite graph-vector flux gives source lower bounds
\[
 \begin{split}
 u\,\dd t&\ge[(1-u)\mathbf1_{(\alpha,\beta)}
                             +q\mathbf1_{(\alpha,\eta)}]\,\dd t,\\
 v\,\dd t&\ge[(1-v)\mathbf1_{(\gamma,\eta)}
                             -p\mathbf1_{(\alpha,\eta)}]\,\dd t.
 \end{split}
\]
For a compact strict corner arc, resolve its increasing-abscissa
displacement $-z'$ in the two wall directions: it assigns first
length $q$ and second length $-p$. Strict contact localization
passes these local bounds from the polygons. The sine/cosine
weighted sum of their right sides is the actual positive graph
projection, namely $2C-T$, which is also the weighted source
total. The nonnegative defects therefore have zero positive
weighted moment and vanish. This proves equality in both bounds.

Put $z_+=N_L(U)(C)$. Since $N_L(U)(-C)=0$, the endpoint
pressures and tangent moments are
\begin{equation}\label{eq:tilted-strip-moments}
 \begin{gathered}
 e_R=\tfrac12+\tfrac h4+\tfrac{3z_+}4,\qquad
 e_L=\tfrac12-\tfrac{3h}4-\tfrac{z_+}4,\\
 T=\int_0^\gamma(1-v)\cos t\,\dd t,
 \quad h=\int_0^\gamma(1-v)\sin t\,\dd t,\\
 T=\int_\beta^L(1-u)\sin t\,\dd t,
 \quad z_+=\int_\beta^L(1-u)\cos t\,\dd t.
 \end{gathered}
\end{equation}
Use the continuous contact energy
\begin{equation}\label{eq:tilted-contact-energy}
 \mathcal H(p,q)=
 \begin{cases}
 (p-1/2)^2+(q+1)^2+3/4,&p,q\ge0,\\
 (p-1)^2+(q+1)^2,&p\le0\le q,\\
 (p-1)^2+(q+1/2)^2+3/4,&p,q\le0.
 \end{cases}
\end{equation}
The formulas agree on the two switching axes. Substitution of
the exact source equations gives
\[
 \mathcal H'=-(q+1)(1-v)\mathbf1_{(0,\gamma)}
                  +(1-p)(1-u)\mathbf1_{(\beta,L)}.
\]
For example in an unclipped core the equations are
$u=(1+q)/2$, $v=(1-p)/2$ and the derivative is zero;
removing floor tangencies replaces the second equation by
$v=-p$ and adds the first displayed term; removing window
tangencies replaces the first by $u=q$ and adds the second.
The same substitution in the outer sign sectors gives the
same terms. Thus every crossing order is included.

Let
\begin{equation}\label{eq:tilted-flux-remainders}
 \begin{gathered}
 I_D=\int_0^\gamma(q+1)(1-v),\qquad
 I_B=\int_\beta^L(1-p)(1-u),\\
 R_D=\int_0^\gamma(1-v(t))\int_0^t u(s)\sin(t-s)\,\dd s\,\dd t,\\
 R_B=\int_\beta^L(1-u(t))\int_t^L v(s)\sin(s-t)\,\dd s\,\dd t.
 \end{gathered}
\end{equation}
Integration variables on the first line are $t$.
Endpoint substitution gives
\[
 \mathcal H(L)-\mathcal H(0)
   =6CT+T^2-(h+z_+)(h+z_+/2).
\]
Variation of constants in the momentum equations gives
\begin{align*}
 q(t)+1&=3C\cos t+(e_R+h)\sin t\\
   &\quad-\int_0^t(1-v(s))\cos(t-s)\,\dd s
         +\int_0^t u(s)\sin(t-s)\,\dd s,\\
 1-p(t)&=(3C+T)\sin t+e_L\cos t\\
   &\quad-\int_t^L(1-u(s))\cos(s-t)\,\dd s
         +\int_t^L v(s)\sin(s-t)\,\dd s.
\end{align*}
For any integrable function $a$, symmetry of the cosine kernel gives
\[
 \int_\xi^\zeta a(t)\int_\xi^t a(s)\cos(t-s)\,\dd s\,\dd t
 =\tfrac12\left[(\int_\xi^\zeta a\cos)^2
                       +(\int_\xi^\zeta a\sin)^2\right].
\]
Applying it to $1-v$ and
$1-u$ yields
\[
 \begin{aligned}
 I_D&=3CT-T^2/2+e_Rh+h^2/2+R_D,\\
 I_B&=3CT+T^2/2+e_Lz_+-z_+^2/2+R_B.
 \end{aligned}
\]
Since the energy change is $I_B-I_D$, the pressures simplify
these identities to
\begin{equation}\label{eq:tilted-two-unit-identity}
 6CT=(z_+-h)(2-h-z_+)/4+R_B-R_D.
\end{equation}
Here $R_D\ge0$, $h+z_+<2$,
$z_+\le\cot\beta\,T\le4T/3$, and
\[
 R_B\le\int_\beta^L(1-u)(1-\sin t)\,\dd t
      \le\frac{1-\sin\beta}{\sin\beta}T\le2T/3.
\]
The first term in \cref{eq:tilted-two-unit-identity} is at most
$z_+/2$. Hence $6CT\le4T/3$, contradicting $C>1/2$, $T>0$.
\end{proof}

\subsection{Paying for a folded companion}

When the companion curvature exceeds one, its tangent can fold back
and the two-unit calculation acquires an error. The common source
occupations bound that error by the unused first-wing curvature. Its
sine moment is already $T$, so the fold can be paid from the same
strip budget.

\begin{lemma}[Folded-wing energy]\label{lem:tilted-fold-energy}
Under the first-unit hypotheses, a companion density exceeding
one also contradicts maximality.
\end{lemma}
\begin{proof}
Let $\alpha<\beta<\tau<\eta$ be the contacts in
\cref{lem:tilted-first-unit-geometry}. That lemma forces $T>0$,
$x_*=-C+T$ and $N_L(U)(-C)=0$.
Use the common occupations $d,\chi$ of
\cref{lem:variation-joint-occupation}. On the core they give
\begin{equation}\label{eq:tilted-fractional-solution}
 u=\frac{b+q\chi}{1+b},\qquad
 v=\frac{d-p\chi}{1+d},\qquad b=\mathbf1_{\{t<\beta\}}.
\end{equation}

Let $\gamma$ be the first zero of $D_y$. We claim $\gamma<\tau$.
There is a zero before $\eta$, since $D_y(0)=-h$ and
$D_y(\eta)=z_y(\eta)>0$.
Suppose $D_y$ has stayed nonpositive until $\tau$.
Then $L_0'=D_y/\sin^2t\le0$ on this initial interval.
At a positive core corner, every earlier second wall is
strictly lower, because
\[
 S_r(z_x(t))=\tan r\,[z_x(t)-L_0(r)]
             <\tan t\,[z_x(t)-L_0(t)]=z_y(t)\quad(r<t).
\]
If the earlier bracket is negative the inequality is immediate;
otherwise both its factors are no greater and the tangent is
strictly smaller. Every later first wall is strictly lower by
global $B_x$ monotonicity. Thus this corner is uniquely exposed
and $\chi=1$ as long as $D_y$ remains nonpositive.
After $\tau$, while this condition persists, $b=d=0$ and
\cref{eq:tilted-fractional-solution} gives
$u=q$, $v=-p$, $p'=q'=-1$. Therefore
$D_y'=(1+p)\sin t<0$. It cannot leave the nonpositive region
before $q$ reaches zero, contradicting $D_y(\eta)>0$.

This reasoning includes $D_y(\tau)=0$ by the following
continuation argument. At $\tau$ the corner is positive and
inside the window, and its two local angular derivatives have
opposite strict signs. The remote-angle gaps are strict and
uniform on compact remote angle sets. Both properties persist
on a neighborhood, giving $\chi=1$ immediately to the right.
Then $p'=-1$; for $p<-1$, the fractional formula gives
$v=(d-p)/(1+d)>1$, forcing $d=0$. Hence $D_y'<0$ there as
well. This excludes an immediate positive departure from zero.
The claim follows, including a possible zero plateau.

Put $S=C+D_x(\gamma)$. At this floor crossing
$D_x(\gamma)=L_0(\gamma)\ge x_*$. The low-wing outer
point has abscissa at most $-C$, so $D_x\le1-C$.
Together with $x_*<1-2C$, these facts give
\begin{equation}\label{eq:tilted-first-floor-moments}
 \begin{gathered}
 0<T\le S\le1,\qquad T\le C,\\
 S=\int_0^\gamma(1-v)\cos t,\qquad
 h=\int_0^\gamma(1-v)\sin t,\\
 T=\int_\beta^L(1-u)\sin t,\qquad
 z_+=\int_\beta^L(1-u)\cos t.
 \end{gathered}
\end{equation}
All four integrals use $\dd t$. The quantity $S$ records the
horizontal displacement up to the first floor crossing; its excess
over $T$ will have a favorable sign in the energy comparison.

On $(\tau,\eta)$, $b=0$ and $p'=q\chi-1-q\le-1$.
Thus $p\le-1$ and $0<q\le1/4$ there. For $w=q-p$,
substitution in the fractional equations gives the exact identity
\[
 w'=(1-\chi)\left[p+q+\frac d{1+d}\right]
             +\frac{\chi d(1+p)}{1+d}\le0.
\]
Both terms are nonpositive since $p+q\le-3/4$ and
$d/(1+d)\le1/2$. As $w(\tau)\le5/4$,
$v\le-p\le5/4-q$ there. Before $\tau$, $v\le1$;
after $\eta$, $v=0$. Hence $v\le5/4$ globally.
Where $v>1$, $d=0$, $v=-p\chi$, $u=q\chi$ and
\begin{equation}\label{eq:tilted-fold-budget}
 (v-1)_+\le1/4-q\le(1-u)/4.
\end{equation}

Use the continuous energy \cref{eq:tilted-contact-energy}.
Before $\beta$, the source equations and $\chi=1$ give
$\mathcal H'=-(q+1)(1-v)(1-d)$.
On the core after $\beta$, direct differentiation gives
\begin{equation}\label{eq:tilted-fractional-energy}
 \mathcal H'-(1-p)(1-u)+(q+1)(1-v)(1-d)
                         =(1-\chi)(p+q).
\end{equation}
The residual is zero before $\tau$ and nonpositive afterwards.
On the final sector $\mathcal H'=1-p$.
On $[0,\gamma]$, $d=0$ and $v\le1$, so its exact contribution
is $-I_D$. All later pieces with $v\le1$ are favorable.
The remaining fold defect is paid by
\cref{eq:tilted-fold-budget}, since $q+1\le5/4$ there:
\[
 \int_{\{v>1\}}(q+1)(v-1)(1-d)
       \le\tfrac5{16}\int_\beta^L(1-u)
       \le\frac{5T}{16\sin\beta}.
\]
It follows that
\begin{equation}\label{eq:tilted-energy-defect}
 \mathcal H(L)-\mathcal H(0)
             \le I_B-I_D+5T/(16\sin\beta).
\end{equation}

The endpoint pressures are those in
\cref{eq:tilted-strip-moments}. The harmonic reconstructions
used in the two-unit proof remain valid, but the first horizontal
moment is now $S$ rather than $T$. Thus
\[
 \begin{split}
 \mathcal H(L)-\mathcal H(0)&=6CT+T^2-(h+z_+)(h+z_+/2),\\
 I_D&=3CS-S^2/2+e_Rh+h^2/2+R_D,\\
 I_B&=3CT+T^2/2+e_Lz_+-z_+^2/2+R_B.
 \end{split}
\]
The remainders have exactly the definitions in
\cref{eq:tilted-flux-remainders}; $R_D\ge0$ because
$v\le1$ before $\gamma$. Substitute these expressions into
\cref{eq:tilted-energy-defect} and use the endpoint pressures.
The result is
\begin{equation}\label{eq:tilted-fold-contradiction}
 \begin{split}
 6CT\le{}&(z_+-h)(2-h-z_+)/4+R_B-R_D\\
 &+(S-T)[(S+T)/2-3C]+5T/(16\sin\beta).
 \end{split}
\end{equation}
The displacement term is nonpositive:
$S\ge T$, $(S+T)/2\le(1+C)/2<3C$.
Using $v\le5/4$ in the backward sine convolution gives
\[
 R_B\le\tfrac54\frac{1-\sin\beta}{\sin\beta}T,
 \qquad z_+\le\cot\beta\,T.
\]
Since $h+z_+<2$, the first term is at most $z_+/2$.
After dropping $-R_D$ and the displacement term, the three remaining
costs are the endpoint niche, the flux remainder, and the fold defect.
Together they bound the right side of
\cref{eq:tilted-fold-contradiction} by
\[
 f(\beta)T,\qquad
 f(\beta)=\frac12\cot\beta
          +\frac54\frac{1-\sin\beta}{\sin\beta}
          +\frac5{16\sin\beta}.
\]
$f$ decreases on $(0,\pi/2]$. Since $\beta>\arctan(3/4)$,
the $3$--$4$--$5$ sine and cosine values give
$f(\beta)<f(\arctan(3/4))<3$.
But $6CT>3T$ for $C>1/2$ and $T>0$, the required contradiction.
\end{proof}
